\documentclass[12pt]{article}
\usepackage[a4paper,margin=1in]{geometry}
\usepackage[utf8]{inputenc}
\usepackage[T1]{fontenc}
\usepackage{amsmath,amssymb}
\usepackage{textcomp}
\usepackage{parskip}
\begin{document}
\begin{center}
{\Large \textbf{2002 STEP I}}
\end{center}

\subsection*{Page 1}
\begin{flushleft}
STEP I, 2002 2\par
\bigskip\textbf{Section A: Pure Mathematics}\par
\bigskip\textbf{Question 1.} Show that the equation of any circle passing through the points of intersection of the ellipse\par
\[(x + 2)2 + 2y2 = 18 and the ellipse 9(x - 1)2 + 16y2 = 25 can be written in the form\]
\[x^2 - 2ax + y^2 = 5 - 4a .\]
\bigskip\textbf{Question 2.} Let f (x) = xm(x - 1)n, where m and n are both integers greater than 1. Show that the\par
curve y = f (x) has a stationary point with 0 < x < 1. By considering f (x), show that this\par
stationary point is a maximum if n is even and a minimum if n is odd.\par
Sketch the graphs of f (x) in the four cases that arise according to the values of m and n.\par
\bigskip\textbf{Question 3.} Show that (a + b)2 6 2a2 + 2b2 .\par
1 1\par
Find the stationary points on the curve y = (a2 cos2 θ + b2 sin2 θ) 2 + (a2 sin2 θ + b2 cos2 θ) 2 ,\par
where a and b are constants. State, with brief reasons, which points are maxima and which\par
are minima. Hence prove that\par
1 1 1\par
\[|a| + |b| 6 (a^2 cos2 θ + b^2 sin2 θ) 2 + (a^2 sin2 θ + b^2 cos2 θ) 2 6 (2a2 + 2b2) 2 .\]
\[4 Give a sketch of the curve y = , for x > 0.\]
\[1 + x^2\]
Find the equation of the line that intersects the curve at x = 0 and is tangent to the curve\par
at some point with x > 0 . Prove that there are no further intersections between the line and\par
the curve. Draw the line on your sketch.\par
\[By considering the area under the curve for 0 6 x 6 1, show that \pi > 3 .\]
Show also, by considering the volume formed by rotating the curve about the y axis, that\par
\[ln 2 > 2/3 .\]
\[1 1 \pi\]
\[[ Note: dx = . ]\]
\[Z 1 + x^2 4\]
\end{flushleft}
\subsection*{Page 2}
\begin{flushleft}
STEP I, 2002 3\par
5 Let\par
\[f (x) = xn + a xn-1 + \cdots + a ,\]
1 n\par
where a , a , . . . , a are given numbers. It is given that f (x) can be written in the form\par
1 2 n\par
\[f (x) = (x + k )(x + k ) \cdots (x + k ) .\]
1 2 n\par
\[By considering f (0), or otherwise, show that k k . . . k = a .\]
1 2 n n\par
Show also that\par
\[(k + 1)(k + 1) \cdots (k + 1) = 1 + a + a + \cdots + a\]
1 2 n 1 2 n\par
and give a corresponding result for (k - 1)(k - 1) \cdots (k - 1).\par
1 2 n\par
Find the roots of the equation\par
\[x^4 + 22x3 + 172x2 + 552x + 576 = 0 ,\]
given that they are all integers.\par
\bigskip\textbf{Question 6.} A pyramid stands on horizontal ground. Its base is an equilateral triangle with sides of\par
length a, the other three sides of the pyramid are of length b and its volume is V . Given that\par
the formula for the volume of any pyramid is 1 \times area of base \times height , show that\par
1 1\par
\[V = a^2(3b2 - a^2) 2 .\]
The pyramid is then placed so that a non-equilateral face lies on the ground. Show that the\par
new height, h, of the pyramid is given by\par
\[a^2(3b2 - a^2)\]
\[h^2 = .\]
\[4b2 - a^2\]
Find, in terms of a and b , the angle between the equilateral triangle and the horizontal.\par
\end{flushleft}
\subsection*{Page 3}
\begin{flushleft}
STEP I, 2002 4\par
7 Let\par
a cos x a sin x\par
\[I = dx and J = dx ,\]
Z sin x + cos x Z sin x + cos x\par
0 0\par
where 0 6 a < 3\pi/4 . By considering I + J and I - J , show that 2I = a + ln(sin a + cos a) .\par
Find also:\par
\[\pi/2 cos x\]
\textbf{(i)} dx , where p and q are positive numbers;\par
Z p sin x + q cos x\par
\[\pi/2 cos x + 4\]
\textbf{(ii)} dx .\par
Z 3 sin x + 4 cos x + 25\par
\bigskip\textbf{Question 8.} I borrow C pounds at interest rate 100α \% per year. The interest is added at the end of each\par
year. Immediately after the interest is added, I make a repayment. The amount I repay at\par
the end of the kth year is R pounds and the amount I owe at the beginning of kth year is\par
k\par
\[C pounds (with C = C). Express C in terms of R (k = 1, 2, . . ., n), α and C and show\]
k 1 n+1 k\par
\[that, if I pay off the loan in N years with repayments given by R = (1 + α)kr , where r is\]
k\par
\[constant, then r = C/N .\]
If instead I pay off the loan in N years with N equal repayments of R pounds, show that\par
R α(1 + α)N\par
\[= ,\]
C (1 + α)N - 1\par
\[and that R/C \approx 27/103 in the case α = 1/50, N = 4 .\]
\end{flushleft}
\subsection*{Page 4}
\begin{flushleft}
STEP I, 2002 5\par
\bigskip\textbf{Section B: Mechanics}\par
2d\par
G\par
P\par
d\par
h\par
α\par
A lorry of weight W stands on a plane inclined at an angle α to the horizontal. Its wheels are\par
a distance 2d apart, and its centre of gravity G is at a distance h from the plane, and halfway\par
between the sides of the lorry. A horizontal force P acts on the lorry through G , as shown.\par
\textbf{(i)} If the normal reactions on the lower and higher wheels of the lorry are equal, show\par
that the sum of the frictional forces between the wheels and the ground is zero.\par
\textbf{(ii)} If P is such that the lorry does not tip over (but the normal reactions on the lower\par
and higher wheels of the lorry need not be equal), show that\par
W tan(α - β) 6 P 6 W tan(α + β) ,\par
\[where tan β = d/h .\]
\end{flushleft}
\subsection*{Page 5}
\begin{flushleft}
STEP I, 2002 6\par
\bigskip\textbf{Question 10.} A bicycle pump consists of a cylinder and a piston. The piston is pushed in with steady\par
speed u. A particle of air moves to and fro between the piston and the end of the cylinder,\par
colliding perfectly elastically with the piston and the end of the cylinder, and always moving\par
parallel with the axis of the cylinder. Initially, the particle is moving towards the piston at\par
speed v. Show that the speed, v , of the particle just after the nth collision with the piston is\par
n\par
\[given by v = v + 2nu.\]
n\par
Let d be the distance between the piston and the end of the cylinder at the nth collision,\par
n\par
and let t be the time between the nth and (n + 1)th collisions. Express d - d in terms\par
n n n+1\par
of u and t , and show that\par
n\par
v + (2n - 1)u\par
\[d = d .\]
n+1 n\par
v + (2n + 1)u\par
Express d in terms of d , u, v and n.\par
n 1\par
d\par
\[In the case v = u, show that ut = .\]
n\par
n(n + 1)\par
P\par
P P\par
1 2 θ\par
φ\par
m km\par
P\par
A particle P of mass m collides with a particle P of mass km which is at rest. No energy is\par
1 2\par
lost in the collision. The direction of motion of P and P after the collision make non-zero\par
1 2\par
angles of θ and φ, respectively, with the direction of motion of P before the collision, as\par
shown. Show that\par
\[sin2 θ + k sin2 φ = k sin2(θ + φ) .\]
Show that, if the angle between the particles after the collision is a right angle, then k = 1 .\par
\end{flushleft}
\subsection*{Page 6}
\begin{flushleft}
STEP I, 2002 7\par
\bigskip\textbf{Section C: Probability and Statistics}\par
12 Harry the Calculating Horse will do any mathematical problem I set him, providing the answer\par
is 1, 2, 3 or 4. When I set him a problem, he places a hoof on a large grid consisting of unit\par
squares and his answer is the number of squares partly covered by his hoof. Harry has circular\par
\[hoofs, of radius 1/4 unit.\]
After many years of collaboration, I suspect that Harry no longer bothers to do the calcu-\par
lations, instead merely placing his hoof on the grid completely at random. I often ask him\par
to divide 4 by 4, but only about 1/4 of his answers are right; I often ask him to add 2 and\par
2, but disappointingly only about \pi/16 of his answers are right. Is this consistent with my\par
suspicions?\par
I decide to investigate further by setting Harry many problems, the answers to which are 1, 2,\par
3, or 4 with equal frequency. If Harry is placing his hoof at random, find the expected value\par
of his answers. The average of Harry's answers turns out to be 2. Should I get a new horse?\par
\bigskip\textbf{Question 13.} The random variable U takes the values +1, 0 and -1 , each with probability 1 . The random\par
variable V takes the values +1 and -1 as follows:\par
\[if U = 1 , then P(V = 1) = 1 and P(V = -1) = 2 ;\]
3 3\par
\[if U = 0 , then P(V = 1) = 1 and P(V = -1) = 1 ;\]
2 2\par
\[if U = -1 , then P(V = 1) = 2 and P(V = -1) = 1 .\]
3 3\par
\textbf{(i)} Show that the probability that both the roots of the equation x2 + U x + V = 0 are\par
real is 1 .\par
\textbf{(ii)} Find the expected value of the larger root of the equation x2 + U x + V = 0 , given\par
that both roots are real.\par
\textbf{(iii)} Find the probability that the roots of the equation\par
\[x^3 + (U - 2V )x^2 + (1 - 2U V )x + U = 0\]
are all positive.\par
\end{flushleft}
\subsection*{Page 7}
\begin{flushleft}
STEP I, 2002 8\par
\bigskip\textbf{Question 14.} In order to get money from a cash dispenser I have to punch in an identification number. I\par
have forgotten my identification number, but I do know that it is equally likely to be any one\par
of the integers 1, 2, . . . , n. I plan to punch in integers in order until I get the right one. I can\par
do this at the rate of r integers per minute. As soon as I punch in the first wrong number,\par
the police will be alerted. The probability that they will arrive within a time t minutes is\par
1 - e\par
-λt,\par
where λ is a positive constant. If I follow my plan, show that the probability of the\par
police arriving before I get my money is\par
n\par
1 - e\par
-λ(k-1)/r\par
.\par
n\par
X\par
\[k=1\]
Simplify the sum.\par
On past experience, I know that I will be so flustered that I will just punch in possible integers\par
at random, without noticing which I have already tried. Show that the probability of the police\par
arriving before I get my money is\par
1 - .\par
n - (n - 1)e-λ/r\par
\end{flushleft}
\end{document}
